% !TeX spellcheck = he_IL
%% ! TEX TS-program = xelatexmk
% above: commands appropriate for TeXShop in mac OS
\documentclass[12pt, a4paper]{article}
\usepackage{etoolbox}
\usepackage{bbm}
%\usepackage{pifont}
\usepackage{tikz}

%%%%%%%%%%%%%%%%%%%%%%%%%%
%%
%%
%%  Exam Template
%%
%%
%%%%%%%%%%%%%%%%%%%%%%%%%%
\newtoggle{testreview}
\toggletrue{testreview}  
\togglefalse{testreview} % uncomment to make it ready for exam unit
\newtoggle{testsolution}
\toggletrue{testsolution}  
\togglefalse{testsolution} % comment to show solutions
\newcounter{version}
\setcounter{version}{2}  % version of multiple-choice

	





\usepackage{graphicx}


\usepackage[a4paper, top=3cm,left=2.5cm, right=2.3cm, bottom=2.7cm]{geometry}
\usepackage{amsmath, amssymb, array, multicol, enumitem}
\usepackage{pdfpages}
\usepackage{comment, graphicx}  
\usepackage{xparse}
\DeclareGraphicsExtensions{.png,.jpg}
\usepackage{pgffor}
\usepackage[colorlinks=true]{hyperref} 
\AtBeginDocument{%
	\hypersetup{pdfborder={0 0 0},
		urlcolor=blue}%
}

\usepackage{wasysym}
\usepackage{fontspec}
\usepackage{xunicode}
\usepackage{polyglossia}
\setdefaultlanguage{hebrew}
\setotherlanguage{english}
%\newfontfamily\hebrewfont[Script=Hebrew, Ligatures=TeX, Scale=1.1, WordSpace =1.05]{David}
%\newfontfamily\hebrewfont[Script=Hebrew, Ligatures=TeX, Scale=1.1, WordSpace=1.05]{David  CLM}
\newfontfamily\hebrewfont[Script=Hebrew]{Times New Roman}
%\newfontfamily\hebrewfont[Script=Hebrew]{Ezra SIL}

% Hebrew monospace font (must support Hebrew)
\newfontfamily\hebrewfonttt[Script=Hebrew]{FreeMono}

%http://tex.stackexchange.com/questions/141231/equation-number-enclosing-parentheses-misplaced-with-polyglossiahebrewamsmath
\makeatletter
%\def\maketag@@@#1{\hbox{\m@th\normalfont\LRE{#1}}}
\def\tagform@#1{\maketag@@@{(\ignorespaces#1\unskip)}}
\makeatother

%
% Hebrew fonts are two small 
%
\fontsize{14}{15}

%
% Package fro proper positioning of logo
%
\usepackage[absolute]{textpos}


% http://tex.stackexchange.com/questions/10589/fix-nested-section-numbers-in-rtl-languages-with-polyglossia/
\makeatletter
\renewcommand{\SepMark}[1]{\def\@SepMark{#1}}
\makeatother
\SepMark{.}
\robustify\hebrewnumeral
\robustify\Hebrewnumeral
\robustify\Hebrewnumeralfinal


% from polyglossia docs
%\gappto\captionshebrew{%
\def\urladdrname{{\itshape כתובת רשת}}
\def\emailaddrname{{\itshape כתובת דואל}}
%}

% https://tex.stackexchange.com/questions/176661/polyglossia-arabic-text-english-bibliography
\makeatletter
\let\orig@MakeUppercase \MakeUppercase
\makeatother
\providecommand{\MakeUppercase}[1]{#1}




%%%%%%%%%%%%%%%%%%%%%%%%%%
%%
%%
%%   User Commands
%%
%%
%%%%%%%%%%%%%%%%%%%%%%%%%%
\newcommand{\heb}[1]{\begin{small}\R{#1}\end{small}}
\newcommand{\hebt}[1]{\begin{scriptsize}\R{#1}\end{scriptsize}}
\newcommand{\cc}{\mathbb{C}}
\newcommand{\nn}{\mathbb{N}}
\newcommand{\zz}{\mathbb{Z}}
%\renewcommand{\phi}{\varphi}
\newcommand{\abs}[1]{\left|#1\right|}
\newcommand{\lr}[1]{\left(#1\right)}


\newcommand{\ff}{\mathbb{F}}
\newcommand{\FF}{\mathbb{F}}
\newcommand{\mtn}{m\times n}
\newcommand{\longr}{\longrightarrow}
\newcommand{\longrl}{\longleftrightarrow}
\newcommand{\ol}{\overline}
\newcommand{\ntn}{n\times n}
\newcommand{\rr}{\mathbb{R}}
\newcommand{\QQ}{\mathbb{Q}}
\newcommand{\RR}{\mathbb{R}}
\newcommand{\im}{\mathrm{Im}}




\newcommand{\twopartdef}[4]
{
	\left\{
		\begin{array}{ll}
			#1 & \mbox{   } #2 \\
			#3 & \mbox{   } #4
		\end{array}
	\right.
}

\newcommand{\solutionword}{\textbf{פתרון:}}
\newcommand{\points}[1]{(#1 נק') }
\iftoggle{testsolution}{%
	\newcommand{\sols}[1]{\solutionword$\;\;$ #1$\square$}
	\newcommand{\spacing}{0cm}
	\newcommand{\questionspacing}{2.5cm}
	\newcommand{\extrapages}{}
}{ %
%\newcommand{\spacing}{12cm}
%\newcommand{\halfspacing}{0.5cm}	
\newcommand{\spacing}{24cm}
\newcommand{\questionspacing}{\spacing}
\newcommand{\sols}[1]{}
\newcommand{\extrapages}{%
	\newpage
	\noindent
	דף זה ריק במכוון ומיועד לשימוש הסטודנט/ית.
	\newpage
	\noindent
	דף זה ריק במכוון ומיועד לשימוש הסטודנט/ית.
}
\iftoggle{testreview}{%
\renewcommand{\spacing}{1cm}
%\newcommand{\questionspacing}{\spacing}
\renewcommand{\extrapages}{}	
}{}	
}






%%%%%%%%%%%%%%%%%%%%%%%%
%%
%%   BEGIN DOCUMENT
%%
%%%%%%%%%%%%%%%%%%%%%%%%

\begin{document}

\selectlanguage{hebrew}


\newcommand{\semester}{%
	סמסטר א' תשפ״ו
}
\newcommand{\course}{%
 אלגברה לינארית להנדסה 
}
\newcommand{\coursenumber}{9321}	
\newcommand{\lecturers}{%
	מרצים: ד״ר י.מייזל, ד״ר נ.גולקו, ד״ר נ.מאיר, ד״ר נ.איידלשטין, ד״ר מ.פורת, ד״ר מ.זיו-אב, ד״ר מ.ברשדסקי
}

\newcommand{\testdate}{{3.2.26}}
\iftoggle{testsolution}{%
	\newcommand{\testtype}{%
		פתרונות למבחן מועד א'
	}
}{%
	\newcommand{\testtype}{%
	\boxed{\text{
	מבחן לדוגמה
	}	}
}}

\newcommand{\duration}{%
	3 שעות.
}
\newcommand{\helpmaterial}{%
מחשבון פשוט בלבד.
}



\thispagestyle{empty}

\begin{textblock*}{9cm}[0.5,0.5](90mm,3cm)
	\includegraphics[width=5cm]{logo.jpg}
\end{textblock*}	

\noindent

\begin{center}
אוניברסיטת בן-גוריון בנגב - היחידה להוראת קורסי יסוד - %
\smallskip
\semester \\
{\large \bf \course - 214.1.\coursenumber}\\
\smallskip
\lecturers\\
\smallskip
\smallskip
\smallskip
{\Large \testtype }\\
\end{center}

\noindent
\textbf{%
תאריך:
} \testdate \quad 
\textbf{משך הבחינה: 
} \duration \quad
\textbf{
חומר עזר: 
}
\helpmaterial \\
דף הנוסחאות מצורף למבחן. אין להכניס כל חומר עזר נוסף.
\subsection*{חלק א': שאלות רב ברירה}
\begin{itemize}
    \item בחלק א' של המבחן יש 8 שאלות רב-ברירה, ומשקל כל שאלה 5 נקודות.
    \item לכל שאלה יש תשובה אחת נכונה בדיוק. יש לסמן תשובה אחת לכל שאלה \textbf{בטבלה בדף זה}, אין צורך בנימוק.
    \item שימו לב להנחיות מילוי הטבלה בעמוד הבא.
\end{itemize}

\begin{center}
\begin{tabular}{|@{}c@{}||@{}c@{}||@{}c@{}||@{}c@{}||@{}c@{}||@{}c@{}||@{}c@{}||@{}c@{}|}
\hline
1 & 2 &  3 & 4 & 5 & 6 & 7 & 8 \\
\hline
\begin{tabular}{cccc}
\\\textcircled{א}\\\\
\textcircled{ב}\\\\
\textcircled{ג}\\\\
\textcircled{ד}\\\\
\end{tabular}
&
\begin{tabular}{cccc}
\\\textcircled{א}\\\\
\textcircled{ב}\\\\
\textcircled{ג}\\\\
\textcircled{ד}\\\\
\end{tabular}

&
\begin{tabular}{cccc}
\\\textcircled{א}\\\\
\textcircled{ב}\\\\
\textcircled{ג}\\\\
\textcircled{ד}\\\\
\end{tabular}
&
\begin{tabular}{cccc}
\\\textcircled{א}\\\\
\textcircled{ב}\\\\
\textcircled{ג}\\\\
\textcircled{ד}\\\\
\end{tabular}
&
\begin{tabular}{cccc}
\\\textcircled{א}\\\\
\textcircled{ב}\\\\
\textcircled{ג}\\\\
\textcircled{ד}\\\\
\end{tabular}
&
\begin{tabular}{cccc}
\\\textcircled{א}\\\\
\textcircled{ב}\\\\
\textcircled{ג}\\\\
\textcircled{ד}\\\\
\end{tabular}
&
\begin{tabular}{cccc}
\\\textcircled{א}\\\\
\textcircled{ב}\\\\
\textcircled{ג}\\\\
\textcircled{ד}\\\\
\end{tabular}
&
\begin{tabular}{cccc}
\\\textcircled{א}\\\\
\textcircled{ב}\\\\
\textcircled{ג}\\\\
\textcircled{ד}\\\\
\end{tabular}
\\
\hline 
\begin{tikzpicture}
  \draw[dotted] (0,0) circle (0.8cm);
\end{tikzpicture}
& 
 \begin{tikzpicture}
  \draw[dotted] (0,0) circle (0.8cm);
\end{tikzpicture}
 &
 \begin{tikzpicture}
  \draw[dotted] (0,0) circle (0.8cm);
\end{tikzpicture}
 &
 \begin{tikzpicture}
  \draw[dotted] (0,0) circle (0.8cm);
\end{tikzpicture}
 &
 \begin{tikzpicture}
  \draw[dotted] (0,0) circle (0.8cm);
\end{tikzpicture}
 &
 \begin{tikzpicture}
  \draw[dotted] (0,0) circle (0.8cm);
\end{tikzpicture}
 &
 \begin{tikzpicture}
  \draw[dotted] (0,0) circle (0.8cm);
\end{tikzpicture}
 &
 \begin{tikzpicture}
  \draw[dotted] (0,0) circle (0.8cm);
\end{tikzpicture}
   \\
 \hline
\end{tabular}
\end{center}


\subsection*{חלק ב': שאלות פתוחות}
\begin{itemize}
    \item יש לבחור בדיוק 3 מתוך 4 שאלות פתוחות. משקל כל שאלה 20 נקודות.

    \item יש לסמן את השאלות שנבחרו בטבלה בדף זה. שאלות שלא סומנו לא תבדקנה!
    \item יש לנמק את התשובות בחלק ב' ולכתוב את כל החישובים הרלוונטים.
    \item יש לענות על כל השאלות בחלק זה על טופס המבחן במקום המיועד אחרי כל שאלה.   
    
    \underline{
    ניתן גם לכתוב מצדו השני של הדף. 
    }
    מחברת הטיוטה לא תבדק!
\end{itemize}
\bigskip

 בחלק ב' יש לבדוק את השאלות הבאות. סמנו את התשובות המתאימות --  3 שאלות לכל היותר!


\begin{center}
\setlength{\arrayrulewidth}{0.6pt}
\renewcommand{\arraystretch}{1.4}
\begin{tabular}{|c|c|c|c|}
\hline
\textcircled{1} & \textcircled{2} & \textcircled{3} & \textcircled{4}\\
\hline
\end{tabular}
\end{center}

\begin{center}
\huge{\textbf{בהצלחה!}}
\end{center}
\pagebreak
\subsection*{הוראות למילוי הטבלה}
\begin{itemize}
    \item[$\bullet$] יש להשתמש בעט שחור או כחול. 
    \item[$\bullet$] יש לסמן תשובה אחת בלבד ע"י השחרת העיגול המתאים.
    \item[$\bullet$]
    בשאלה בה לא יושחר אף עיגול,
    הניקוד יהיה 0.
    \item[$\bullet$] בשאלה בה יושחרו יותר מעיגול אחד, הניקוד יהיה 0.
.
    \item[$\bullet$] כדי לתקן סימון, יש לסמן X
    על התשובה הלא נכונה ולהשחיר את התשובה הנכונה.

    \item[$\bullet$]
    \textbf{שימו לב:}
    לא ניתן לבחור סעיף שסומן עליו X
    לכן הקפידו לסמן את תשובתכם כשאתם בטוחים בה.
\end{itemize}

\subsection*{דוגמאות}


\begin{center}
\begin{tabular}{|c||c||c||c||c||c||c||c||c||c|}
\hline
1 & 2 &  3 & 4 & 5 & 6 & 7 & 8 & 9 & 10 \\
\hline
\begin{tabular}{c}
\\\Large{\CIRCLE}\\\\
\ \  \hspace{-8pt}\raisebox{0pt}[0pt][0pt]{\Large\CIRCLE}%
\  \hspace{4pt}\llap{\raisebox{-4pt}[0pt][0pt]{ $\scalebox{3}{$\times$}$}}\hspace{-8pt}
\\\\
\textcircled{ג}\\\\
\textcircled{ד}\\\\
\end{tabular}
&
\begin{tabular}{cccc}
\\\textcircled{א}\\\\
\Large{\CIRCLE}\\\\
\textcircled{ג}\\\\
\textcircled{ד}\\\\
\end{tabular}
&
\begin{tabular}{cccc}
\\
\textcircled{א}
\\\\
\ \  \hspace{-8pt}\raisebox{0pt}[0pt][0pt]{\Large\CIRCLE}%
 \ \hspace{4pt}\llap{\raisebox{-4pt}[0pt][0pt]{ $\scalebox{3}{$\times$}$}}\hspace{-8pt}
\\\\
\Large{\CIRCLE}\\\\
\ \  \hspace{-8pt}\raisebox{0pt}[0pt][0pt]{\Large\CIRCLE}%
 \ \hspace{4pt}\llap{\raisebox{-4pt}[0pt][0pt]{ $\scalebox{3}{$\times$}$}}\hspace{-8pt}
\\\\
\end{tabular}
&
\begin{tabular}{cccc}
\\\textcircled{א}\\\\
\textcircled{ב}\\\\
\textcircled{ג}\\\\
\textcircled{ד}\\\\
\end{tabular}
&
\begin{tabular}{cccc}
\\\hspace{-8pt}\ \  \raisebox{0pt}[0pt][0pt]{\Large\CIRCLE}%
 \ \hspace{4pt}\llap{\raisebox{-4pt}[0pt][0pt]{ $\scalebox{3}{$\times$}$}}\hspace{-8pt}
\\\\
\textcircled{ב}\\\\
\textcircled{ג}\\\\
\textcircled{ד}\\\\
\end{tabular}
&
\begin{tabular}{cccc}
\\\hspace{-8pt}\ \  \raisebox{0pt}[0pt][0pt]{\Large\CIRCLE}%
 \ \hspace{4pt}\llap{\raisebox{-4pt}[0pt][0pt]{ $\scalebox{3}{$\times$}$}}\hspace{-8pt}
\\\\
\hspace{-8pt}\ \  \raisebox{0pt}[0pt][0pt]{\Large\CIRCLE}%
 \ \hspace{4pt}\llap{\raisebox{-4pt}[0pt][0pt]{ $\scalebox{3}{$\times$}$}}\hspace{-8pt}
\\\\
\textcircled{ג}\\\\
\textcircled{ד}\\\\
\end{tabular}
&
\begin{tabular}{cccc}
\\\Large{\CIRCLE}\\\\
\Large{\CIRCLE}\\\\
\textcircled{ג}\\\\
\textcircled{ד}\\\\
\end{tabular}
&
\begin{tabular}{cccc}
\\\Large{\CIRCLE}\\\\
\Large{\CIRCLE}\\\\
\Large{\CIRCLE}\\\\
\textcircled{ד}\\\\
\end{tabular}
&
\begin{tabular}{cccc}
\\\hspace{-8pt}\ \  \raisebox{0pt}[0pt][0pt]{\Large\CIRCLE}%
 \ \hspace{4pt}\llap{\raisebox{-4pt}[0pt][0pt]{ $\scalebox{3}{$\times$}$}}\hspace{-8pt}
\\\\
\Large{\CIRCLE}\\\\
\textcircled{ג}\\\\
\Large{\CIRCLE}\\\\
\end{tabular}
&
\begin{tabular}{cccc}
\\\hspace{-8pt}\ \  \raisebox{0pt}[0pt][0pt]{\Large\CIRCLE}%
 \ \hspace{4pt}\llap{\raisebox{-4pt}[0pt][0pt]{ $\scalebox{3}{$\times$}$}}\hspace{-8pt}
\\\\
\Large{\CIRCLE}
\\\\
\hspace{-8pt}\ \  \raisebox{0pt}[0pt][0pt]{\Large\CIRCLE}%
 \ \hspace{4pt}\llap{\raisebox{-4pt}[0pt][0pt]{ $\scalebox{3}{$\times$}$}}\hspace{-8pt}
\\\\
\Large{\CIRCLE}\\\\
\end{tabular}
\\
\hline
\end{tabular}
\ 

\ 

\begin{itemize}[label=$\bullet$]
    \item שאלות 1,2,3 עם סימונים תקינים עבור תשובות:
    \begin{enumerate}
        \item א
        \item ב
        \item ג
    \end{enumerate}
    \item בשאלות 4,5,6 חסר סימון והניקוד יהיה 0.
    \item בשאלות 7,8,9,10
    יש יותר מסימון אחד והניקוד יהיה 0
\end{itemize}


\end{center}


%%%
%%%
%%%   End of Header and Instructions
%%%
%%%

\newpage

\underline{{\Large
}}


\begin{enumerate}

%%%%%%%%%%%%%%%%%%%%%%%%%%%%%%%%%%%%%%
%%  Problems
%%%%%%%%%%%%%%%%%%%%%%%%%%%%%%%%%%%%%%


%%%%% שאלה 2 %%%%%%%

\section*{
חלק א' 
(40 נק')
להשחיר את התשובות בטבלה!}
\item\points{5}
 יהי $V$ מרחב וקטורי ויהיו $U_1,U_2,U_3 \subseteq V$ תתי-מרחבים.
בחרו את הטענה הנכונה:

\begin{enumerate}
\item
אם $U_1\cap U_2=U_1\cap U_3$ אז $U_2=U_3$.

\item
אם $\dim(U_1)+\dim(U_2)-\dim(U_3)>\dim(V)$ אז $U_1+U_2+U_3=V$.

\item
אם $\dim(U_1+U_2+U_3)=\dim(U_1+U_2)$ אז $U_3\subseteq U_1+U_2$.

\item
אם $U_1\nsubseteq U_2$ אז $U_1\cap U_2=\{0\}$.
\end{enumerate}

\noindent\rule{\textwidth-30pt}{0.4pt}

\item\points{5}
ידוע כי $(1,1,1),(1,1,2),(1,2,3)$ הם פתרונות של ממ״ל $AX=b$,

וידוע כי $(2,2,2)$ אינו פתרון של הממ״ל $AX=b$. בחרו את הטענה הנכונה:

\begin{enumerate}
\renewcommand{\labelenumi}{\hebrew{enumi}.}
\item
$(0,1,1)$ פתרון של $AX=b$.

\item
לכל $a\in\mathbb{R}$ מתקיים: $(1,1,a)$ פתרון של $AX=b$.

\item
$(0,0,0)$ פתרון של $AX=b$.

\item
לכל $a\in\mathbb{R}$ מתקיים: $(a,a,a)$ פתרון של $AX=b$.
\end{enumerate}

\noindent\rule{\textwidth-30pt}{0.4pt}

%%%%% שאלה 2 %%%%%%%

\item\points{5}
 יהיו $B,C$ בסיסים של $\mathbb{R}^2$ עבורם $a,b\in\mathbb{R}$ מקיימים
\[
[I]^B_C=\frac{1}{\sqrt{2}}
\begin{bmatrix}
a & 1\\
1 & -1
\end{bmatrix}
\qquad
[I]^C_B=\frac{1}{\sqrt{2}}
\begin{bmatrix}
b & 1\\
1 & -1
\end{bmatrix}
\]
נזכיר כי $[I]^B_C$ היא מטריצת המעבר מבסיס $B$ ל-$C$. בחרו את הטענה הנכונה:

\begin{tabular}{@{}p{0.48\linewidth} p{0.48\linewidth}@{}}
(א) $a=2,b=1$ & (ג) $b=-1,a=-1$\\
(ב) $b=1,a=0$ & (ד) $b=1,a=1$\\
\end{tabular}

\noindent\rule{\textwidth-30pt}{0.4pt}


\item\points{5}
יהי $a,b\in\mathbb{R}$. נתון כי הקבוצה
\[
\left\{
\begin{pmatrix}1\\[2pt]1\\[2pt]a\end{pmatrix},
\begin{pmatrix}1\\[2pt]-2\\[2pt]a\end{pmatrix},
\begin{pmatrix}-1\\[2pt]b\\[2pt]1\end{pmatrix}
\right\}
\]
היא בסיס אורתוגונלי של המרחב הווקטורי $\mathbb{R}^3$.
בחרו את הטענה הנכונה:

\begin{tabular}{@{}p{0.48\linewidth} p{0.48\linewidth}@{}}
(א) $a=-1,b=2$ & (ג) $a=1,b=0$\\
(ב) $a=-1,b=-1$ & (ד) $a=1,b=1$\\
\end{tabular}

\noindent\rule{\textwidth-30pt}{0.4pt}

\item\points{5}
יהיו
\[
A=
\begin{pmatrix}
i & 1\\
1 & -i
\end{pmatrix},
\qquad
B=
\begin{pmatrix}
-i & 0\\
2 & i
\end{pmatrix}.
\]
נסמן $c=\det(A^2+2AB+B^2)$. בחרו את הטענה הנכונה:

\begin{tabular}{@{}p{0.48\linewidth} p{0.48\linewidth}@{}}
(א) $c=-7$ & (ג) $c=-3$\\
(ב) $c=9$  & (ד) $c=25$\\
\end{tabular}



%%%
%%%%%%
%%%

\noindent\rule{\textwidth-30pt}{0.4pt}



%%%
%%%%%%
%%%
\item\points{5}
 נתון כי $A$ לכסינה וכן הפולינום האופייני של   $A$ הוא
\[
p_A(x)=x^5-5x^3+x^2.
\]
בחרו את הטענה הנכונה:

\begin{tabular}{@{}p{0.48\linewidth} p{0.48\linewidth}@{}}
(א) $\mathrm{rank}(A)=3$ & (ג) $\mathrm{rank}(A)=5$\\
(ב) $\mathrm{rank}(A)=2$ & (ד) $\mathrm{rank}(A)=4$\\
\end{tabular}

\noindent\rule{\textwidth-30pt}{0.4pt}


%%%
%%%%%%
%%%

\item\points{5}
תהי $T:\mathbb{R}^3\to\mathbb{R}^2$ העתקה לינארית המקיימת
\[
T\!\begin{pmatrix}1\\0\\1\end{pmatrix}=\begin{pmatrix}1\\2\end{pmatrix},
\qquad
T\!\begin{pmatrix}1\\2\\3\end{pmatrix}=\begin{pmatrix}7\\-4\end{pmatrix},
\qquad
T\!\begin{pmatrix}1\\1\\1\end{pmatrix}=\begin{pmatrix}2\\0\end{pmatrix}.
\]
בחרו את הטענה הנכונה:

\begin{tabular}{@{}p{0.48\linewidth} p{0.48\linewidth}@{}}
\[
T\!\begin{pmatrix}3\\-1\\1\end{pmatrix}=\begin{pmatrix}-2\\10\end{pmatrix}
\] (א)
&
\[
T\!\begin{pmatrix}3\\-1\\1\end{pmatrix}=\begin{pmatrix}1\\1\end{pmatrix}
\] (ב)
\\[10pt]
\[
T\!\begin{pmatrix}3\\-1\\1\end{pmatrix}=\begin{pmatrix}5\\12\end{pmatrix}
\] (ג)
&
\[
T\!\begin{pmatrix}3\\-1\\1\end{pmatrix}=\begin{pmatrix}3\\-6\end{pmatrix}
\] (ד)
\end{tabular}

\noindent\rule{\textwidth-30pt}{0.4pt} 

\item \points{5}
נתון מישור $H$ שעובר דרך הנקודות $(0,-2,1),(1,1,-1),(2,1,0)$.
מי מבין המישורים הבאים מקביל ל-$H$? בחרו את התשובה הנכונה:

\begin{tabular}{@{}p{0.48\linewidth} p{0.48\linewidth}@{}}
(א) $\{(x,y,z)\in\mathbb{R}^3 \mid x+y-z=3\}$ &
(ג) $\{(x,y,z)\in\mathbb{R}^3 \mid 2x-2y-2z=1\}$\\[6pt]
(ב) $\{(x,y,z)\in\mathbb{R}^3 \mid 2x-y+2z=-1\}$ &
(ד) $\{(x,y,z)\in\mathbb{R}^3 \mid x+y-z=-3\}$\\
\end{tabular}

\noindent\rule{\textwidth-30pt}{0.4pt}

\end{enumerate}

%%%
%%%%%%
%%%

 \pagebreak

\section*{חלק ב' (60 נק') - יש לבחור 3 מתוך 4 ולסמן בעמוד הראשון.}
יש לענות במקומות המתאימים בעמודים הבאים.
\begin{enumerate}
\newcommand{\firstQ}{
\item \points{20}
תהי
\[A = \begin{pmatrix}
    -1 & -1 & 4\\
    -1 & 0 & 3\\
    0 & -1 & 1
\end{pmatrix} \]
\begin{enumerate}
    \item\points{10} האם $A$ לכסינה מעל $\rr$?
    האם $A$ לכסינה מעל $\cc$?
    \item\points{5} מצאו מטריצה $B$ כך ש-$B^{3}=A$.
    \item\points{5} מצאו את כל הערכים העצמיים של המטריצה
    $2A^3-3A^5$.
\end{enumerate}

}
\firstQ
\noindent\rule{\textwidth-30pt}{0.4pt}

\newcommand{\secondQ}{
\item\points{20}
יהיו $a,b\in\mathbb{R}$, ותהי $T:\mathbb{R}^3\to\mathbb{R}^3$ העתקה לינארית המקיימת
\[
T\!\begin{pmatrix}1\\1\\1\end{pmatrix}=\begin{pmatrix}1\\1\\0\end{pmatrix},
\qquad
T\!\begin{pmatrix}1\\1\\0\end{pmatrix}=\begin{pmatrix}a\\1\\0\end{pmatrix},
\qquad
T\!\begin{pmatrix}1\\0\\0\end{pmatrix}=\begin{pmatrix}1\\1\\b\end{pmatrix}.
\]

\begin{enumerate}
\item[(א)] (5 נק') מצאו את המטריצה $A=[T]^E_E$ כאשר $E$ הוא הבסיס הסטנדרטי של $\mathbb{R}^3$ וחשבו את $\det(A)$.
\item[(ב)] (5 נק') מצאו עבור כל $a,b\in\mathbb{R}$ בסיס ל-$\operatorname{Im}(T)$.
\item[(ג)] (10 נק') מצאו לכל $0\le k\le 2$ את כל הערכים $a,b\in\mathbb{R}$ עבורם $\dim(\ker(T))=k$ ומצאו בסיס של $\ker(T)$.
\end{enumerate}
}
\secondQ
\noindent\rule{\textwidth-30pt}{0.4pt}
\newcommand{\thirdQ}{
\item 
תהי
\[
A=
\begin{pmatrix}
1 & 3 & 4\\
1 & 1 & 2\\
-1 & 1 & 0
\end{pmatrix}
\in M_{3\times 3}(\mathbb{R}).
\]
יהיו $W=N(A)$ מרחב הפתרונות של הממ״ל $AX=0$, ויהי $U=\operatorname{Col}(A)$ מרחב העמודות של $A$.
\begin{enumerate}
\item[(א)] (8 נק') מצאו בסיס ל-$U$ ובסיס ל-$W$.
\item[(ב)] (8 נק') מצאו בסיס אורתונורמלי ל-$U\cap W$ ובסיס אורתונורמלי ל-$U+W$.
\item[(ג)] (4 נק') מצאו מטריצה $B\in M_{3\times 3}(\mathbb{R})$ כך שמתקיים $U=N(B)$.
\end{enumerate}
}
\thirdQ
\noindent\rule{\textwidth-30pt}{0.4pt}
\newcommand{\forthQ}{
    \item \points{20}
    הסעיפים הבאים לא קשורים אחד לשני:
    \begin{enumerate}
        \item\points{10} תהי
        $T:\rr^3\to\rr^3$ ויהיו 
        $v_1,v_2\in \rr^3$ לא קו-לינאריים כך שמתקיים
        $T(v_1) = 3T(v_2)$.

        הוכיחו כי
        $T$ לא על.
        \item\points{10} תהי $S:\rr^3\to\rr^2$.
        הוכיחו כי ישנם שני וקטורים
        $v_1,v_2\in\rr^3$ שאינם קו-לינאריים,
        
        כך שמתקיים
        $S(v_1) = S(v_2)$.
    \end{enumerate}
}\forthQ
\pagebreak
\setcounter{enumi}{0}
    \firstQ
\noindent\rule{\textwidth-30pt}{0.4pt}
\underline{פיתרון:}
    
    \pagebreak \ 
    \ 
    
    \pagebreak

    
    \secondQ
    \noindent\rule{\textwidth-30pt}{0.4pt}
    \underline{פיתרון:}


    \pagebreak \ \pagebreak

    \thirdQ
    \noindent\rule{\textwidth-30pt}{0.4pt}
    \underline{פיתרון:}


    \pagebreak \ \pagebreak

    \forthQ
    \noindent\rule{\textwidth-30pt}{0.4pt}
    \underline{פיתרון:}


    \pagebreak \ \pagebreak
\end{enumerate}

    \pagebreak 
    
    עמוד נוסף לכתיבת המשך פיתרונות
    
    \pagebreak
    עמוד נוסף לכתיבת המשך פיתרונות

 \pagebreak

 % --- דף הנוסחאות ---
\section*{דף הנוסחאות}

\begin{enumerate}
\item $M_{m\times n}(F)$ זו קבוצת המטריצות בגודל $n\times m$
($m$ שורות, $n$ עמודות)
עם רכיבים מ-$F$.

\noindent\rule{\textwidth-30pt}{0.4pt}

\item לכל מטריצה $A$, נסמן על ידי $(A)_{i,j}$ את האיבר של $A$ הנמצא בשורה ה-$i$ ובעמודה ה-$j$.

\noindent\rule{\textwidth-30pt}{0.4pt}

\item פיתוח דטרמיננטה של $A\in M_{n\times n}(F)$ לפי השורה ה-$i$ או לפי העמודה ה-$j$:
\[
\det(A)=\sum_{i=1}^{n}(-1)^{i+j}a_{ij}\det\!\big(M_{ij}(A)\big)
=\sum_{j=1}^{n}(-1)^{i+j}a_{ij}\det\!\big(M_{ij}(A)\big)
\]
כאן $M_{ij}(A)$ זה המינור ה-$i,j$ של המטריצה $A$.

\noindent\rule{\textwidth-30pt}{0.4pt}

\item אם $T:V\to W$ העתקה לינארית, אם $B=\{v_1,\ldots,v_n\}$ בסיס סדור של $V$ ואם $C$ בסיס סדור של $W$,
אז המטריצה המייצגת של $T$ ביחס לבסיסים $B$ ו-$C$ מוגדרת על ידי
\[
[T]^B_C=
\begin{pmatrix}
\vline & & \vline\\
[T(v_1)]_C & \cdots & [T(v_n)]_C\\
\vline & & \vline
\end{pmatrix}.
\]

\noindent\rule{\textwidth-30pt}{0.4pt}

\item לכל $u=(a_1,\ldots,a_n),\; w=(b_1,\ldots,b_n)\in\mathbb{R}^n$

נסמן את המכפלה הפנימית הסטנדרטית ב-$\mathbb{R}^n$ (המכפלה הסקלרית):
\[
 u\cdot w=a_1\cdot b_1+\cdots+a_n\cdot b_n.
\]

\noindent\rule{\textwidth-30pt}{0.4pt}

% --- דף הנוסחאות ---


\item תהליך גראם-שמיט (גרסה א):
אם $W$ תת-מרחב של $\rr^n$, ואם $\{v_1,\ldots,v_m\}$ בסיס של $W$

אז $\{w_1,\ldots,w_m\}$ בסיס אורתוגונלי של $W$ כאשר
\[
w_1=v_1,\qquad
w_k=v_k-\sum_{i=1}^{k-1}\frac{ v_k\cdot w_i}{ w_i\cdot w_i}\,w_i
\quad (k=2,\ldots,m).
\]
\item תהליך גראם-שמיט (גרסה ב):
בהנתן סדרה בת"ל
$\{v_1,\ldots,v_k\}\subseteq \rr^n$, נבנה סדרה אורתונורמלית באופן הבא:
\[
w_1=\frac{v_1}{\|v_1\|},\qquad
w_i'=v_i-\sum_{l=1}^{i-1} (v_i\cdot w_l)\,w_l,\qquad w_i = \frac{w_i'}{\|w_i\|}
\]

\noindent\rule{\textwidth-30pt}{0.4pt}

\item השורשים של המשוואה הריבועית
$
0=c+bx+ax^2,\qquad a\ne 0,
$
נתונים על ידי
\[
x_{1,2}=\frac{-b\pm\sqrt{b^2-4ac}}{2a}.
\]

\noindent\rule{\textwidth-30pt}{0.4pt}

\item מטריצה $A\in M_{n\times n}(\ff)$ דומה למטריצה $B\in M_{n\times n}(\ff)$ 

אם
קיימת 
מטריצה הפיכה
$P\in M_{n\times n}(\ff)$ כך ש-
$A=PBP^{-1}$.
\end{enumerate}

\end{document}